
\subsubsection{3.1 Fidelity and Richness Framework}

Feedback quality is decomposed into two dimensions:

\textbf{Signal Fidelity}: The probability that feedback is factually correct. Execution feedback achieves near-perfect fidelity (tests are deterministic). AI feedback has lower fidelity (approximately 39\% correctness based on Self-Refine analysis).

\textbf{Semantic Richness}: The amount of explanatory content about why code is wrong. Execution feedback has minimal richness (only pass/fail). AI feedback provides detailed explanations.

These dimensions are orthogonal. A method achieving high fidelity and high richness requires filtering rich signals through a fidelity check.

\subsubsection{3.2 EVAF Mechanism}

EVAF treats execution as a verification layer for AI feedback:

\begin{enumerate}
\item \textbf{AI Critique Generation}: Given problem description and failing code, the feedback model generates critique and proposed fix using a structured prompt.
\item \textbf{Code Extraction}: Fenced code blocks (``\texttt{python ... }``) are parsed from the AI response using regular expression matching.
\item \textbf{Execution Gating}: The extracted fix is executed against unit tests in a subprocess with 3.0 second timeout and network isolation.
\item \textbf{Accept/Reject Decision}: If all tests pass, the suggestion is accepted; otherwise, it is rejected with classified error type.
\end{enumerate}

\subsubsection{3.3 Implementation Architecture}

Seven Python modules implement the pipeline:

\begin{itemize}
\item \textbf{config.py}: Configuration constants including model identifiers, generation parameters, and paths.
\item \textbf{data.py}: HumanEval loading via HuggingFace datasets (\texttt{openai\textbackslash \_humaneval}).
\item \textbf{model.py}: \texttt{BaselineModel} wrapping CodeT5-770M (\texttt{Salesforce/codet5-large}) and \texttt{FeedbackModel} wrapping CodeLlama-7b-Instruct (\texttt{codellama/CodeLlama-7b-Instruct-hf}, loaded in float16).
\item \textbf{gating.py}: Code extraction using regex, error classification, and sandboxed subprocess test execution.
\item \textbf{metrics.py}: Accept rate, coverage, and rejection breakdown computation.
\item \textbf{visualize.py}: Figure generation (gate metrics bar chart, accept distribution histogram, rejection breakdown pie chart).
\item \textbf{train.py}: Pipeline orchestration combining all modules.
\end{itemize}

\subsubsection{3.4 Accept Rate as Viability Indicator}

\begin{table}[htbp]
\centering
\resizebox{\columnwidth}{!}{%
\begin{tabular}{ll}
\toprule
Accept Rate & Interpretation \\
\midrule
<10\% & EVAF degenerates to pure execution feedback; AI nearly always wrong \\
20-60\% & Selective filtering; AI provides substantive verified feedback \\
>90\% & Gating unnecessary; AI nearly always correct \\
\bottomrule
\end{tabular}
}
\end{table}
